% Part: first-order-logic
% Chapter: axiomatic-deduction
% Section: soundness

\documentclass[../../../include/open-logic-section]{subfiles}

\begin{document}

\iftag{FOL}
      {\olfileid{fol}{axd}{sou}}
      {\olfileid{pl}{axd}{sou}}
      
\olsection{Kasahihan}

\begin{explain}
Sistem !!^a{derivation}, kayata deduksi aksiomatik, iku \emph{sah} yen
ora bisa !!{derive} bab kang sejatine ora lumaku. Mula kasahihan iku
sawijining sipat kaslametan kang dijamin kanggo sistem
!!{derivation}. Gumantung marang sipat ing teori bukti kang dirembug,
kita kepengin ngerti, upamane, yen
\begin{enumerate}
\item saben $!A$ kang !!{derivable} iku valid;
\item yen $!A$ iku !!{derivable} saka sawetara formula liyane~$\Gamma$,
  formula iku uga dadi akibate;
\item yen sawijining himpunan !!{formula}s~$\Gamma$ ora konsisten,
  himpunan iku ora bisa dicukupi.
\end{enumerate}
Iki sipat-sipat wigati saka sistem !!a{derivation}. Yen salah sijine
ora lumaku, sistem !!{derivation} kasebut cacad---sistem iku bakal
!!{derive} kakehan. Mula, netepake kasahihan sistem
!!a{derivation} iku wigati banget.
\end{explain}

\begin{prop}
  Yen $!A$ minangka aksioma, mula
  \iftag{FOL}
    {$\Sat{M}{!A}[s]$ kanggo saben !!{structure}~$\Struct{M}$ lan penetapan~$s$.}
    {$\pSat{v}{!A}$ kanggo saben !!{valuation}~$\pAssign{v}$.}
\end{prop}

\begin{proof}
\iftag{FOL}{Kita kudu mriksa yen kabeh aksioma iku valid. Tuladhane,
  iki kasus kanggo \olref[qua]{ax:q1}: upamana $t$ iku !!{free for}
  $x$ ing $!A$, lan upamana
  $\Sat{M}{\lforall[x][!A]}[s]$. Banjur miturut definisi panyukupan,
  kanggo saben $\varAssign{s'}{s}{x}$, uga $\Sat{M}{!A}[s']$, lan
  mligine iki lumaku nalika $s'(x) = \Value{t}{M}[s]$. Miturut
  \olref[syn][ext]{prop:ext-formulas},
  $\Sat{M}{\Subst{!A}{t}{x}}[s]$. Iki nuduhake yen
  $\Sat{M}{(\lforall[x][!A] \lif \Subst{!A}{t}{x})}[s]$.}{Gawea tabel
  kabeneran kanggo saben aksioma supaya bisa dipriksa yen kabeh iku
  tautologi.}
\end{proof}

\begin{thm}[Kasahihan]
\ollabel{thm:soundness}
Kanggo himpunan premis lan kesimpulan kang kabeh awujud ukara, yen
$\Gamma \Proves !A$, mula $\Gamma \Entails !A$. Cathetan koreksi
sumber OLPL-037: watesan ukara iki digawe cetha amarga semantik
tataran kapisan lan kasus aturan kuantor ing bukti nggunakake watesan
kasebut.
\end{thm}

\begin{proof}
Kanthi induksi miturut cacah aplikasi aturan inferensi ing
!!{derivation} kang formula pungkasane $!A$ saka $\Gamma$. Cathetan koreksi sumber
OLPL-038: ukuran induksi iki cocog karo basis lan langkah induktif ing
sumber. Yen ora ana langkah kang dhasare aturan inferensi, mula
kabeh !!{formula}s ing !!{derivation} kasebut minangka conto aksioma
utawa ana ing~$\Gamma$. Miturut proposisi sadurunge, kabeh aksioma iku
\iftag{FOL}{valid}{tautologi}, mula yen $!A$ iku aksioma, banjur
$\Gamma \Entails !A$. Yen $!A \in \Gamma$, mula kanthi langsung
$\Gamma \Entails !A$.

Yen langkah pungkasan !!{derivation} kang formula pungkasane~$!A$
nduweni dhasar modus ponens, mula ana !!{formula}s $!B$ lan $!B \lif !A$ ing
!!{derivation}, lan hipotesis induksi lumaku kanggo perangan
!!{derivation} kang mungkasi ing !!{formula}s kasebut (amarga
perangan kasebut ngemot paling ora siji langkah kang dhasare aturan
inferensi luwih sithik). Dadi, miturut hipotesis induksi, $\Gamma
\Entails !B$ lan $\Gamma \Entails !B \lif !A$. Banjur $\Gamma
\Entails !A$ amarga saben struktur utawa valuasi kang nyukupi Gamma uga
nyukupi rong formula kasebut; iki uga ndherek saka Teorema Deduksi Semantis lan
transitivitas akibat, yaiku
\iftag{FOL}
      {\olref[syn][sem]{thm:sem-deduction}}
      {\olref[pl][syn][sem]{thm:sem-deduction}}.

\iftag{FOL}{Saiki upamana langkah pungkasan nduweni dhasar
\QR. Banjur langkah kasebut awujud $!C \lif
\lforall[x][!B(x)]$, lan ana langkah sadurunge $!C \lif !B(c)$ kanthi
$c$ ora muncul ing $\Gamma$, $!C$, utawa $\lforall[x][!B(x)]$. Miturut
hipotesis induksi, $\Gamma \Entails !C \lif !B(c)$. Miturut
\olref[syn][sem]{thm:sem-deduction}, $\Gamma \cup \{!C\} \Entails
!B(c)$. Cathetan koreksi sumber OLPL-039: tandha metavariabel formula
ing rong formula kuantor kasebut dibalekake supaya cocog karo kabeh
panggonan sabanjure ing bukti iki.

Gatekna sawijining struktur~$\Struct{M}$ kanthi
$\Sat{M}{\Gamma \cup \{!C\}}$. Kita kudu nuduhake yen
$\Sat{M}{\lforall[x][!B(x)]}$. Amarga
$\lforall[x][!B(x)]$ iku !!a{sentence}, iki ateges kita kudu nuduhake
yen kanggo saben penetapan variabel~$s$, $\Sat{M}{!B(x)}[s]$
(\olref[syn][ass]{prop:sat-quant}). Amarga $\Gamma \cup \{!C\}$ mung
dumadi saka ukara, $\Sat{M}{!D}[s]$ kanggo kabeh $!D \in \Gamma$
miturut \olref[syn][sat]{defn:satisfaction}. Tetepna $\Struct{M'}$ padha
karo $\Struct{M}$ kajaba $\Assign{c}{M'} = s(x)$. Amarga $c$ ora
muncul ing~$\Gamma$ utawa~$!C$, $\Sat{M'}{\Gamma \cup \{!C\}}$
miturut \olref[syn][ext]{cor:extensionality-sent}. Amarga $\Gamma
\cup \{!C\} \Entails !B(c)$, $\Sat{M'}{!B(c)}$. Amarga $!B(c)$ iku
!!a{sentence}, $\Sat{M'}{!B(c)}[s]$ miturut
\olref[syn][ass]{prop:sentence-sat-true}. $\Sat{M'}{!B(x)}[s]$ yen lan
mung yen $\Sat{M'}{!B(c)}[s]$ miturut
\olref[syn][ext]{prop:ext-formulas} (elinga yen $!B(c)$ mung
$\Subst{!B(x)}{c}{x}$). Dadi, $\Sat{M'}{!B(x)}[s]$. Amarga $c$ ora
muncul ing~$!B(x)$, miturut
\olref[syn][ext]{prop:extensionality}, $\Sat{M}{!B(x)}[s]$. Nanging
$s$ iku penetapan variabel sawenang-wenang, mula
$\Sat{M}{\lforall[x][!B(x)]}$. Dadi, $\Gamma \cup \{!C\} \Entails
\lforall[x][!B(x)]$. Miturut \olref[syn][sem]{thm:sem-deduction},
$\Gamma \Entails !C \lif \lforall[x][!B(x)]$. Cathetan koreksi sumber
OLPL-039 uga lumaku kanggo argumen formula ing pratelan panyukupan kang
wis dibalekake ing paragraf iki.

Kasus nalika $!A$ nduweni dhasar \QR{} nanging awujud
$\lexists[x][!B(x)] \lif !C$ ditinggalake minangka gladhen.}{}
\end{proof}

\tagprob{FOL}
\begin{prob}
  Rampungna bukti \olref[fol][axd][sou]{thm:soundness}.
\end{prob}
\tagendprob

\begin{cor}
\ollabel{cor:weak-soundness}
Yen $\Proves !A$, mula $!A$ iku \iftag{FOL}{valid}{tautologi}.
\end{cor}

\begin{cor}
\ollabel{cor:consistency-soundness}
Yen $\Gamma$ bisa dicukupi, mula himpunan kasebut konsisten.
\end{cor}

\begin{proof}
Kita mbuktekake kontraposisine. Upamana $\Gamma$ ora konsisten. Banjur
$\Gamma \Proves \lfalse$, yaiku ana !!a{derivation} kang formula
pungkasane $\lfalse$ saka~$\Gamma$. Miturut \olref{thm:soundness}, saben
\iftag{FOL}{!!{structure}~$\Struct{M}$}{!!{valuation}~$\pAssign{v}$}
kang nyukupi $\Gamma$ mesthi nyukupi~$\lfalse$. Amarga
\iftag{FOL}{$\Sat/{M}{\lfalse}$ kanggo saben
  !!{structure}~$\Struct{M}$}{$\pSat/{v}{\lfalse}$ kanggo saben
  !!{valuation}~$\pAssign{v}$}, ora ana
\iftag{FOL}{$\Struct{M}$}{$\pAssign{v}$} kang bisa nyukupi $\Gamma$,
yaiku $\Gamma$ ora bisa dicukupi.
\end{proof}


\end{document}
